% Ground truth for photo-p1.jpg + photo-p2.jpg (simulated photos), unattended run, edited level.
% The last line of page 2 is cut off in the photo.
% Only what the scorer needs: content blocks, display equations and marks.
\section{Oscillateur harmonique}

\begin{definition}
\src{p01}
Un système de position $x(t)$ est un \emph{oscillateur harmonique} s'il vérifie
\begin{equation}\label{eq-oh}
  \ddot{x} + \omega_0^2\, x = 0
\end{equation}
où $\omega_0 > 0$ est la pulsation propre.
\end{definition}

\begin{property}
\src{p01}
Les solutions de \eqref{eq-oh} sont les fonctions $x(t) = A\cos(\omega_0 t + \varphi)$,
avec $A \geq 0$ et $\varphi \in \R$.
\end{property}

\begin{remark}
\src{p01}
\struck{On admet que} Le mouvement est périodique, de période
\[
  \suspect{T_0 = \frac{\pi}{\omega_0}}{x(t) = A cos(omega_0 t + phi) a pour période 2 pi / omega_0 (propriété p01)}.
\]
\end{remark}

\src{p01}
\textbf{Énergie.} Pour une masse $m$ reliée à un ressort de raideur $k$, on a
$\omega_0 = \sqrt{k/m}$ et
\[
  E = \frac{1}{2} m \dot{x}^2 + \frac{1}{2} k x^2 .
\]

\begin{example}
\src{p02}
Pour $m = \qty{0,2}{\kilogram}$ et $k = \qty{20}{\newton\per\metre}$, on obtient
$\omega_0 = \qty{10}{\radian\per\second}$.
\end{example}

\begin{algorithm}[htbp]
\caption{Euler symplectique\src{p02}}
\begin{algorithmic}
  \For{$n = 0$ \textbf{à} $N-1$}
    \State $v_{n+1} \gets v_n - h\,\omega_0^2\, x_n$
    \State $x_{n+1} \gets x_n + h\, v_{n+1}$
  \EndFor
\end{algorithmic}
\end{algorithm}

\begin{exercise}
\src{p02}
Montrer que $E$ est constante au cours du mouvement.
\end{exercise}

\src{p02}
\missing{dernière ligne coupée sur la photo (commence par « Indication »)}
